\clearpage
\VerdierPhysicalPageAnchor{0216}%
\begingroup\small\setlength{\parskip}{0pt}\setlength{\emergencystretch}{1.5em}
\setlength{\abovedisplayskip}{4pt plus 2pt minus 2pt}%
\setlength{\belowdisplayskip}{4pt plus 2pt minus 2pt}%

\VerdierSemanticAnchor{ch3d-3-2-8-cont}
où $Z''_{n-1}$ est la somme amalgamée :

\VerdierSemanticAnchor{ch3d-3-2-8-pushout-square}
\begin{center}
\begin{tikzcd}[column sep=4.2em,row sep=2.8em]
Y \arrow[r,"u_n"] \arrow[d,"g"']
  & Z_{n-1} \arrow[d,"g''"] \\
Y'' \arrow[r,"u''_n"'] & Z''_{n-1}
\end{tikzcd}
\end{center}
et où $u''_{n-1}$ est l'unique morphisme défini par les égalités
$u''_{n-1}g''=u_{n-1}$ et $u''_{n-1}u''_n=0$.

\VerdierSemanticAnchor{ch3d-3-2-9}
\medskip\noindent\textbf{Proposition 3.2.9.}\quad
\textit{a) On a les égalités :}
\[
\delta^n(S^n(X,Y)*f)=\delta^n(S^n(X,Y))f,
\]
\[
\delta^n(g*S^n(X,Y))=g[n]\,\delta^n(S^n(X,Y)).
\]

\VerdierSemanticAnchor{ch3d-3-2-9-b-diagram}
\noindent\textit{b) Tout diagramme commutatif :}
\begin{center}
\resizebox{\textwidth}{!}{%
$\begin{tikzcd}[ampersand replacement=\&,column sep=1.75em,row sep=2.7em]
S^n(X,Y)=0 \arrow[r]
  \& Y \arrow[r] \arrow[d,"g"']
  \& Z_{n-1} \arrow[r] \arrow[d]
  \& \cdots \arrow[r]
  \& Z_1 \arrow[r] \arrow[d]
  \& Z_0 \arrow[r] \arrow[d]
  \& X \arrow[r] \arrow[d,"f"]
  \& 0 \\
{S'}^n(X',Y')=0 \arrow[r]
  \& Y' \arrow[r]
  \& Z'_{n-1} \arrow[r]
  \& \cdots \arrow[r]
  \& Z'_1 \arrow[r]
  \& Z'_0 \arrow[r]
  \& X' \arrow[r]
  \& 0
\end{tikzcd}$%
}
\end{center}
\noindent\textit{fournit l'égalité :}

\VerdierSemanticAnchor{ch3d-3-2-9-1}
\begin{equation}
\delta^n(g*S^n(X,Y))=\delta^n({S'}^n(X',Y')*f).
\tag{3.2.9.1}
\end{equation}

\VerdierSemanticAnchor{ch3d-3-2-9-proof}
Lorsque $n=1$, ces propriétés résultent des propriétés du
$\delta$-foncteur canonique
$\operatorname{comp}(\mathcal A)\longrightarrow D(\mathcal A)$ (1.3).
Le cas général s'obtient par récurrence sur $n$, en utilisant (3.2.5).

\VerdierSemanticAnchor{ch3d-3-2-10}
\medskip\noindent\textbf{3.2.10.}\quad
Soient $S^n(X,Y)$ et ${S'}^n(X',Y')$ deux éléments appartenant
respectivement aux ensembles $E_{\mathcal A}^{n}(X,Y)$ et
$E_{\mathcal A}^{n}(X',Y')$ et considérons le morphisme :

\VerdierSemanticAnchor{ch3d-3-2-10-direct-sum-map}
\[
\delta^n\bigl(S^n(X,Y)\oplus {S'}^n(X',Y')\bigr)
\colon X\oplus X'\longrightarrow (Y\oplus Y')[n].
\]
On a l'égalité :

\VerdierSemanticAnchor{ch3d-3-2-10-1}
\begin{equation}
\delta^n\bigl(S^n(X,Y)\oplus {S'}^n(X',Y')\bigr)
=
\begin{pmatrix}
\delta^n(S^n(X,Y))&0\\
0&\delta^n({S'}^n(X',Y'))
\end{pmatrix},
\tag{3.2.10.1}
\end{equation}
ainsi qu'il résulte immédiatement de la définition des applications
$\delta_n$.
\endgroup
